
%%%%%%%%%%%%%%%%%%%%%%%%%%%%%%%%%%%%%%%%%%%%%%%%%%%%%%%%%%%%%
\section{模型的改进与推广}

\subsection{模型的改进}

（1）\textbf{振打同步性约束。}引入同步因子$S=(\sum\delta_i)^2/\sum\delta_i^2$衡量多电场同时振打的风险，$S=4$为完全同步（峰值最大），$S=1$为完全错峰。数据中$T_1/T_2\approx1.003$、$T_3/T_4\approx1.001$，前两电场和后两电场各自同步振打。在优化中增加错峰约束$T_i\neq T_j$或相位错开，降低瞬时排放峰值。

（2）\textbf{贝叶斯参数估计完善。}当前Laplace近似因限幅数据残差非正态，95\% CI覆盖率仅3\%。可改用Bootstrap重采样或MCMC获得更可靠的参数后验分布，在论文中报告参数的置信范围而非仅点估计。

（3）\textbf{振打过频系数$r$标定。}$r=0.5$为工程先验，可结合振打瞬态实验数据或CFD模拟标定$r$值，量化过频振打的二次扬尘效应。